\section*{Bab \ref{ch:b2:basins-compared} --- Perbandingan Cekungan Sedimen}

\begin{solution}{exo:b2:basins-compared:1}
(a) Intrakratonik: penurunan lambat yang berdenyut selama ratusan juta tahun, tanpa pensesaran
besar; penyebabnya diperdebatkan. (b) Celah: penipisan mekanis kerak selama peregangan,
lalu kontraksi termal selagi litosfer mendingin. (c) Tepi pasif: penurunan termal
yang diwarisi dari tahap celah, diperbesar oleh beban baji sedimen yang tebal. (d) dan (e)
Depan-pegunungan: litosfer yang dilengkungkan ke bawah oleh beban sabuk sesar dorong, yaitu Alpen dan Himalaya.
\end{solution}

\begin{solution}{exo:b2:basins-compared:2}
(a) $Y = 3000\times(3330 - 2300)/(3330 - 1030) = 3000\times1030/2300 = \qty{1343}{m}$.
(b) $1343 + 200 = \qty{1543}{m}$. (c) Tanpa bebannya, batuan dasar akan terletak pada \qty{1343}{m}:
$1 - 1343/3000 = \qty{55}{\percent}$ tebal itu diberi ruang oleh beban.
\end{solution}

\begin{solution}{exo:b2:basins-compared:3}
(a) $\beta = 35/25 = 1.4$. (b) $125/1.4 = \qty{89.3}{km}$. (c) $S_i = 4334\times(1 - 1/1.4) =
\qty{1238}{m}$.
\end{solution}

\begin{solution}{exo:b2:basins-compared:4}
(a) $\tau = (1.25\times10^{5})^2/(\pi^2\times8\times10^{-7}) = \qty{1.979e15}{s}$, yaitu
$1.979\times10^{15}/3.156\times10^{13} = \qty{62.7}{Ma}$. (b) $62.7\times0.8 = \qty{50.2}{Ma}$.
(c) $62.7\times(100/125)^2 = \qty{40.1}{Ma}$. (d) $1 - e^{-1} = \qty{63}{\percent}$; $1 -
e^{-3} = \qty{95}{\percent}$.
\end{solution}

\begin{solution}{exo:b2:basins-compared:5}
(a) $0.63\,e^{-1000/1960} = 0.378$; $0.63\,e^{-2000/1960} = 0.227$. (b) $T_s = 1000 -
0.63\times1960\times(0.6004 - 0.3605) = \qty{704}{m}$. (c) $m = 2700\times704 + 1030\times296 =
\qty{2.206e6}{kg.m^{-2}}$; $\bar\rho = \qty{2206}{kg.m^{-3}}$.
\end{solution}

\begin{solution}{exo:b2:basins-compared:6}
(a) $z_2' \leftarrow 704 + 1234.8\,(1 - e^{-z_2'/1960})$: 1077, 1226, 1278, 1296, 1301, 1303,
1304; lapisan itu setebal \qty{1304}{m}. (b) $(2700\times704 + 1030\times600)/1304 =
\qty{1932}{kg.m^{-3}}$. (c) Sebesar \qty{304}{m}, atau \qty{30}{\percent}: air yang telah diperas keluar oleh
pemadatan.
\end{solution}

\begin{solution}{exo:b2:basins-compared:7}
(a) $1 - e^{-100/62.7} = 0.7971$. Untuk $\beta = 1.5$: $3207\times0.4135\times0.7971 =
\qty{1057}{m}$; untuk $\beta = 2$: $3207\times0.6366\times0.7971 = \qty{1627}{m}$. (b) Ya: 1100 sampai
\qty{1400}{m} terletak di antaranya. (c) $\frac{\beta}{\pi}\sin\frac{\pi}{\beta} = 1250/(3207\times0.7971)
= 0.4890$, sehingga $\beta = 1.6 + 0.1\times(0.4890 - 0.4705)/(0.5204 - 0.4705) = 1.64$, di dalam rentang
peregangan 50 sampai \qty{100}{\percent} yang disimpulkan untuk cekungan itu.
\end{solution}

\begin{solution}{exo:b2:basins-compared:8}
(a) Massa jenis isian menggantikan $\rho_w$ dalam penyebut. (b) Pembilangnya adalah
$125\,000\times(1/3)\times74.69 = \num{3.112e6}$; $S_i = 3.112\times10^{6}/(3184.4 - 2300) =
\qty{3519}{m}$. (c) Di bawah air, $3.112\times10^{6}/2154.4 = \qty{1445}{m}$: cekungan yang terisi sedimen
2.44 kali lebih dalam.
\end{solution}

\begin{solution}{exo:b2:basins-compared:9}
(a) $2171 - 100\times3330/2300 = 2171 - 145 = \qty{2026}{m}$. (b) $Y$ berubah meter demi meter
bersama $W_d$: dari 1971 sampai \qty{2371}{m}. (c) Kedalaman air, $\pm\qty{200}{m}$, dibandingkan dengan
\qty{145}{m} untuk muka laut.
\end{solution}

\begin{solution}{exo:b2:basins-compared:10}
(a) $S_i$ bertanda sama dengan kurung dalam pembilangnya, karena penyebutnya, $3184.4 - 1030$,
positif; nilainya negatif, yaitu pengangkatan, bila kurung itu negatif. (b) $(\rho_m - \rho_c)t_c/a =
\alpha T_1\rho_m/2$ memberikan $t_c = a\,\alpha T_1\rho_m/(2(\rho_m - \rho_c)) = 125\times0.04372\times
3330/1060 = \qty{17.17}{km}$. (c) $t_c = 17.17/(1 - 0.04372\times17.17/250) = \qty{17.22}{km}$.
(d) Di tempat kerak lebih tipis daripada sekitar \qty{17}{km} sebelum peregangan: di litosfer samudra,
atau di kerak benua yang sudah ditipiskan oleh pembentukan celah sebelumnya.
\end{solution}

\begin{solution}{exo:b2:basins-compared:11}
(a) Sebelum peregangan, kolom itu memuat litosfer bermassa $M = \rho_c t_c(1 - \alpha T_1
t_c/(2a)) + \rho_m(a - t_c)(1 - \alpha T_1(a + t_c)/(2a))$; sesudahnya, $\rho_w S_i + M/\beta +
\rho_m(1 - \alpha T_1)(a - S_i - a/\beta)$. Dengan menyamakan keduanya, $S_i[\rho_m(1 - \alpha T_1) - \rho_w] =
(1 - 1/\beta)[\rho_m(1 - \alpha T_1)a - M]$, dan dengan menjabarkan $M$, $\rho_m(1 - \alpha T_1)a - M =
(\rho_m - \rho_c)t_c(1 - \alpha T_1 t_c/(2a)) - \rho_m\alpha T_1 a/2$: itulah rumusnya. (b) Profil akhirnya adalah $T_1 z/a$; profil awal melebihinya
sebesar $T_1 z(\beta - 1)/a$ di atas $a/\beta$ dan sebesar $T_1(1 - z/a)$ di bawahnya. Integral-integralnya adalah $T_1
a(\beta - 1)/(2\beta^2)$ dan $T_1 a(\beta - 1)^2/(2\beta^2)$, yang jumlahnya $T_1 a(1 - 1/\beta)/2$.
Penyusutan sebesar $\alpha$ kali nilai ini, yang dibebani air, memberikan ungkapan itu. (c) $3.28\times10^{-5}
\times3330\times1333\times125\,000\times0.5/(2\times2300) = \qty{1978}{m}$, dibandingkan dengan $3207\times
2/\pi = \qty{2042}{m}$: suku pertama melebih-lebihkan penurunan akhir sebesar \qty{64}{m},
\qty{3}{\percent}.
\end{solution}

\begin{solution}{exo:b2:basins-compared:12}
(a) $3207\times0.6366\times(1 - e^{-150/62.7}) = 2042\times0.9086 = \qty{1855}{m}$. (b) $1855\times
2300/(3330 - 2400) = \qty{4588}{m}$. (c) 6.9 sampai \qty{13.8}{km}. (d) Model itu mengabaikan beberapa
hal yang dapat memperdalam cekungan: litosfer mantel yang teregang lebih kuat daripada kerak, tanggapan regional
(fleksural) terhadap beban sedimen, pengangkatan sementara tepi itu selama pembentukan celah
yang disusul penurunan, atau pasokan sedimen yang begitu besar sehingga cekungan terisi berlebih dan terbebani
melampaui muka laut di paparan.
\end{solution}

\begin{solution}{pb:b2:basins-compared:1}
\textbf{1.} $0.49\,e^{-4420/3704} = 0.149$ di puncak; $0.49\,e^{-6030/3704} = 0.096$ di dasar.
\textbf{2.} $T_s = 1610 - 0.49\times3704\times(0.3032 - 0.1963) = \qty{1416}{m}$.
\textbf{3.} Serpih: $1110 - 0.63\times1960\times(e^{-3310/1960} - e^{-4420/1960}) = 1110 - 1234.8
\times(0.1847 - 0.1049) = \qty{1011}{m}$.
\textbf{4.} Pasir lempungan: $2680\times1004 + 1030\times696 = \qty{3.408e6}{kg.m^{-2}}$; kapur:
$2710\times1409 + 1030\times201 = \qty{4.025e6}{kg.m^{-2}}$; serpih: $2700\times1011 + 1030\times99 =
\qty{2.832e6}{kg.m^{-2}}$; pasir: $2650\times1416 + 1030\times194 = \qty{3.952e6}{kg.m^{-2}}$. Totalnya
\qty{1.4217e7}{kg.m^{-2}}, $\bar\rho_s = 14\,216\,910/6030 = \qty{2358}{kg.m^{-3}}$.
\textbf{5.} $z_2' \leftarrow 1416 + 1815\,(1 - e^{-z_2'/3704})$: 1993, 2171, 2221, 2235, 2238, 2239,
2239: \qty{2239}{m}.
\textbf{6.} $(2650\times1416 + 1030\times823)/2239 = \qty{2055}{kg.m^{-3}}$.
\textbf{7.} $Y = 2239\times(3330 - 2055)/2300 + 200 = 1241 + 200 = \qty{1441}{m}$.
\textbf{8.} $Y = 6030\times(3330 - 2358)/2300 + 100 = 2548 + 100 = \qty{2648}{m}$. Satuan-satuan itu
diukur pada kedalamannya kini: satuan-satuan itu sudah sepadat keadaannya kini.
\textbf{9.} $(2171 - 1441)/50 = \qty{14.6}{m}$ per juta tahun dari 150 sampai \qty{100}{Ma};
$(2453 - 2171)/40 = \qty{7.05}{m}$ dari 100 sampai \qty{60}{Ma}; $(2648 - 2453)/60 = \qty{3.25}{m}$
sejak itu: penurunannya melambat.
\textbf{10.} $6030 + 100 = \qty{6130}{m}$, 2.31 kali penurunan tektonik.
\textbf{11.} $(6130 - 2648)/6130 = \qty{57}{\percent}$.
\textbf{12.} $2171 - 100\times3330/2300 = \qty{2026}{m}$.
\textbf{13.} Titik pada akhir pembentukan celah, \qty{150}{Ma}: $S_i = \qty{1441}{m}$, jika tidak ada penurunan
termal yang terjadi selama \qty{15}{Ma} pembentukan celah.
\textbf{14.} $\beta = 1/(1 - 1441/4334) = 1/0.6675 = 1.50$.
\textbf{15.} Kerak setebal $35/1.50 = \qty{23.3}{km}$ dan litosfer setebal $125/1.50 = \qty{83.3}{km}$.
\textbf{16.} $(1.50/\pi)\sin(\pi/1.50) = 0.4775\times0.8660 = 0.4135$; $3207\times0.4135 =
\qty{1326}{m}$.
\textbf{17.} $S_i = 4334/3 = \qty{1445}{m}$. Pada \qty{100}{Ma}, $t = \qty{50}{Ma}$: $1326\times(1 -
e^{-50/62.7}) = \qty{729}{m}$, $Y = \qty{2173}{m}$; pada \qty{60}{Ma}: $1326\times0.7620 =
\qty{1010}{m}$, $Y = \qty{2455}{m}$; kini: $1326\times0.9086 = \qty{1205}{m}$, $Y = \qty{2650}{m}$.
\textbf{18.} Model itu \qty{2}{m} lebih dalam di setiap titik: 2173 dibandingkan dengan 2171, 2455 dibandingkan dengan 2453,
2650 dibandingkan dengan 2648.
\textbf{19.} Kedalaman air purba yang hanya diketahui dengan ketelitian seratus meter atau lebih; muka laut masa lalu
yang berbeda dari muka laut kini; pembentukan celah yang berlangsung \qty{15}{Ma}, yang selama itu sebagian kalor
sudah hilang; fleksur alih-alih isostasi lokal; mantel yang teregang lebih kuat atau lebih lemah daripada kerak;
kerak yang tebalnya bukan \qty{35}{km} sebelum pembentukan celah.
\textbf{20.} $1326\times(e^{-150/62.7} - e^{-200/62.7}) = \qty{67}{m}$.
\textbf{21.} $66.6\times2300/(3330 - 2400) = \qty{165}{m}$.
\textbf{22.} $1 - e^{-150/62.7} = \qty{91}{\percent}$.
\textbf{23.} $\tau\ln2 = \qty{43.5}{Ma}$ setelah akhir pembentukan celah: pada \qty{106.5}{Ma}.
\textbf{24.} $\beta = \textbf{1.50}$: kerak teregang menjadi satu setengah kali lebarnya.
\end{solution}
